\documentclass[a4paper,12pt]{article}
\usepackage[a4paper, margin=1.8cm]{geometry}
\usepackage{tikz}
\usetikzlibrary{calc}
\usepackage{amsmath}
\usepackage{amssymb}
\usepackage{array}
\usepackage{xcolor}
\usepackage{mdframed}
\usepackage{enumitem}
\usepackage[T1]{fontenc}
\usepackage{textcomp}

% ---- Colours ----
\definecolor{slice1}{RGB}{70,130,180}
\definecolor{slice2}{RGB}{240,140,40}
\definecolor{slice3}{RGB}{70,170,70}
\definecolor{slice4}{RGB}{210,70,70}
\definecolor{slice5}{RGB}{150,90,190}
\definecolor{lightgrey}{RGB}{245,245,245}
\definecolor{boxblue}{RGB}{220,235,250}
\definecolor{boxgreen}{RGB}{220,245,220}
\definecolor{boxorange}{RGB}{255,240,220}

\newcommand{\piechart}[2]{%
  \begin{tikzpicture}
    \pgfmathsetmacro{\r}{#1}%
    \pgfmathsetmacro{\startA}{90}%
    \foreach \lbl/\ang/\col in {#2} {%
      \pgfmathsetmacro{\endA}{\startA - \ang}%
      \fill[\col!25] (0,0) -- (\startA:\r cm) arc(\startA:\endA:\r cm) -- cycle;%
      \draw[black, line width=1.2pt] (0,0) -- (\startA:\r cm);%
      \draw[black, line width=1.2pt] (0,0) -- (\endA:\r cm);%
      \pgfmathsetmacro{\midA}{(\startA + \endA)/2}%
      \pgfmathsetmacro{\lR}{\r * 0.65}%
      \node[font=\small\bfseries, align=center] at (\midA:\lR cm) {\lbl};%
      \xdef\startA{\endA}%
    }%
    \draw[black, line width=1.2pt] (0,0) circle (\r cm);%
    \fill[black] (0,0) circle (2.5pt);%
  \end{tikzpicture}%
}

\pagestyle{empty}
\parindent=0pt

\begin{document}

\begin{center}
  {\LARGE\bfseries Worked Solutions: Drawing Pie Charts with a Protractor}\\[6pt]
  {\large Teacher Copy --- Completed tables and angle workings}
\end{center}

\vspace{4pt}

\begin{mdframed}[backgroundcolor=boxgreen, linecolor=slice3,
                 linewidth=1.5pt, roundcorner=6pt]
  \textbf{Key formula:}\quad
  $\text{sector angle}=\dfrac{\text{amount}}{\text{total}}\times 360^\circ$
\end{mdframed}

\vspace{6pt}

% ===========================
%  CHART 1
% ===========================
\begin{mdframed}[backgroundcolor=boxblue, linecolor=slice1,
                 linewidth=1.2pt, roundcorner=6pt]
\textbf{Chart 1 -- Favourite Pets (class of 36 students)}

\medskip
\begin{minipage}[c]{0.42\textwidth}
  \centering
  \piechart{4.0}{Dog/120/slice1,Cat/90/slice2,Fish/80/slice3,Rabbit/50/slice4,Other/20/slice5}
\end{minipage}%
\hfill
\begin{minipage}[c]{0.54\textwidth}
  \centering
  \small
  \renewcommand{\arraystretch}{1.5}
  \begin{tabular}{|l|c|c|c|}
    \hline
    \textbf{Pet} & \textbf{No.} & \textbf{Fraction} & \textbf{Angle} \\
    \hline
    Dog    & 12 & $\frac{1}{3}$ & $120^\circ$ \\
    \hline
    Cat    &  9 & $\frac{1}{4}$ & $90^\circ$ \\
    \hline
    Fish   &  8 & $\frac{2}{9}$ & $80^\circ$ \\
    \hline
    Rabbit &  5 & $\frac{5}{36}$ & $50^\circ$ \\
    \hline
    Other  &  2 & $\frac{1}{18}$ & $20^\circ$ \\
    \hline
    \textbf{Total} & \textbf{36} & \textbf{1} & \textbf{360\textdegree} \\
    \hline
  \end{tabular}
\end{minipage}

\medskip
\textbf{Angle working:}
\[
\frac{12}{36}\times360=120^\circ,\quad
\frac{9}{36}\times360=90^\circ,\quad
\frac{8}{36}\times360=80^\circ,
\]
\[
\frac{5}{36}\times360=50^\circ,\quad
\frac{2}{36}\times360=20^\circ.
\]
Check: \(120+90+80+50+20=360^\circ\).

\textbf{Answers:}
\begin{enumerate}[leftmargin=1.6em, itemsep=2pt]
  \item Completed chart shown above. The remaining sectors are Fish \(80^\circ\), Rabbit \(50^\circ\), Other \(20^\circ\).
  \item Dog, because \(\frac{12}{36}=\frac13\).
  \item Cat: 9 students corresponds to a \(90^\circ\) sector.
\end{enumerate}
\end{mdframed}

\vspace{6pt}

% ===========================
%  CHART 2
% ===========================
\begin{mdframed}[backgroundcolor=boxblue, linecolor=slice2,
                 linewidth=1.2pt, roundcorner=6pt]
\textbf{Chart 2 -- Sports Club Memberships (36 members)}

\medskip
\begin{minipage}[c]{0.42\textwidth}
  \centering
  \piechart{4.0}{Football/140/slice1,Swimming/90/slice2,Tennis/60/slice3,Basketball/50/slice4,Other/20/slice5}
\end{minipage}%
\hfill
\begin{minipage}[c]{0.54\textwidth}
  \centering
  \small
  \renewcommand{\arraystretch}{1.5}
  \begin{tabular}{|l|c|c|c|}
    \hline
    \textbf{Sport} & \textbf{No.} & \textbf{Fraction} & \textbf{Angle} \\
    \hline
    Football   & 14 & $\frac{7}{18}$ & $140^\circ$ \\
    \hline
    Swimming   &  9 & $\frac{1}{4}$ & $90^\circ$ \\
    \hline
    Tennis     &  6 & $\frac{1}{6}$ & $60^\circ$ \\
    \hline
    Basketball &  5 & $\frac{5}{36}$ & $50^\circ$ \\
    \hline
    Other      &  2 & $\frac{1}{18}$ & $20^\circ$ \\
    \hline
    \textbf{Total} & \textbf{36} & \textbf{1} & \textbf{360\textdegree} \\
    \hline
  \end{tabular}
\end{minipage}

\medskip
\textbf{Angle working:}
\[
\frac{14}{36}\times360=\frac{7}{18}\times360=140^\circ,\quad
\frac{9}{36}\times360=90^\circ,\quad
\frac{6}{36}\times360=60^\circ,
\]
\[
\frac{5}{36}\times360=50^\circ,\quad
\frac{2}{36}\times360=20^\circ.
\]
Check: \(140+90+60+50+20=360^\circ\).

\textbf{Answers:}
\begin{enumerate}[leftmargin=1.6em, itemsep=2pt]
  \item Completed chart shown above.
  \item Swimming has a \(90^\circ\) sector because \(\frac{9}{36}=\frac14\).
  \item Tennis fraction \(=\frac{6}{36}=\frac16\).
  \item The three smallest sports are Tennis, Basketball and Other:
  \[
  60^\circ+50^\circ+20^\circ=130^\circ.
  \]
\end{enumerate}
\end{mdframed}

\vspace{6pt}

% ===========================
%  CHART 3
% ===========================
\begin{mdframed}[backgroundcolor=boxblue, linecolor=slice3,
                 linewidth=1.2pt, roundcorner=6pt]
\textbf{Chart 3 -- How Students Spend Most of Their Pocket Money (36 students)}

\medskip
\begin{minipage}[c]{0.42\textwidth}
  \centering
  \piechart{4.0}{Food \& Drink/120/slice1,Games/90/slice2,Clothes/60/slice3,Savings/60/slice4,Other/30/slice5}
\end{minipage}%
\hfill
\begin{minipage}[c]{0.54\textwidth}
  \centering
  \small
  \renewcommand{\arraystretch}{1.5}
  \begin{tabular}{|l|c|c|c|}
    \hline
    \textbf{Category} & \textbf{No.} & \textbf{Fraction} & \textbf{Angle} \\
    \hline
    Food \& Drink & 12 & $\frac{1}{3}$ & $120^\circ$ \\
    \hline
    Games         &  9 & $\frac{1}{4}$ & $90^\circ$ \\
    \hline
    Clothes       &  6 & $\frac{1}{6}$ & $60^\circ$ \\
    \hline
    Savings       &  6 & $\frac{1}{6}$ & $60^\circ$ \\
    \hline
    Other         &  3 & $\frac{1}{12}$ & $30^\circ$ \\
    \hline
    \textbf{Total} & \textbf{36} & \textbf{1} & \textbf{360\textdegree} \\
    \hline
  \end{tabular}
\end{minipage}

\medskip
\textbf{Angle working:}
\[
\frac{12}{36}\times360=120^\circ,\quad
\frac{9}{36}\times360=90^\circ,\quad
\frac{6}{36}\times360=60^\circ,
\]
\[
\frac{6}{36}\times360=60^\circ,\quad
\frac{3}{36}\times360=30^\circ.
\]
Check: \(120+90+60+60+30=360^\circ\).

\textbf{Answers:}
\begin{enumerate}[leftmargin=1.6em, itemsep=2pt]
  \item Completed chart shown above.
  \item Clothes and Savings both have 6 students, so they both have a sector angle of \(60^\circ\).
\end{enumerate}
\end{mdframed}

\vspace{6pt}

% ===========================
%  CHART 4
% ===========================
\begin{mdframed}[backgroundcolor=boxblue, linecolor=slice4,
                 linewidth=1.2pt, roundcorner=6pt]
\textbf{Chart 4 -- Colours of Cars in a Car Park (60 cars)}

\medskip
\begin{minipage}[c]{0.42\textwidth}
  \centering
  \piechart{4.0}{White/120/slice1,Black/90/slice2,Silver/72/slice3,Red/54/slice4,Other/24/slice5}
\end{minipage}%
\hfill
\begin{minipage}[c]{0.54\textwidth}
  \centering
  \small
  \renewcommand{\arraystretch}{1.6}
  \begin{tabular}{|l|c|c|c|}
    \hline
    \textbf{Colour} & \textbf{No.} & \textbf{Fraction} & \textbf{Angle} \\
    \hline
    White  & 20 & $\frac{1}{3}$ & $120^\circ$ \\
    \hline
    Black  & 15 & $\frac{1}{4}$ & $90^\circ$ \\
    \hline
    Silver & 12 & $\frac{1}{5}$ & $72^\circ$ \\
    \hline
    Red    &  9 & $\frac{3}{20}$ & $54^\circ$ \\
    \hline
    Other  &  4 & $\frac{1}{15}$ & $24^\circ$ \\
    \hline
    \textbf{Total} & \textbf{60} & \textbf{1} & \textbf{360\textdegree} \\
    \hline
  \end{tabular}
\end{minipage}

\medskip
\textbf{Table working:}
\begin{align*}
\text{White: } \frac{20}{60}=\frac13 &\Rightarrow \frac13\times360=120^\circ,\\
\text{Black: } \frac{15}{60}=\frac14 &\Rightarrow \frac14\times360=90^\circ,\\
\text{Silver: } \frac{12}{60}=\frac15 &\Rightarrow \frac15\times360=72^\circ,\\
\text{Red: } \frac{9}{60}=\frac{3}{20} &\Rightarrow \frac{3}{20}\times360=54^\circ,\\
\text{Other: } \frac{4}{60}=\frac{1}{15} &\Rightarrow \frac{1}{15}\times360=24^\circ.
\end{align*}

\textbf{Answers:}
\begin{enumerate}[leftmargin=1.6em, itemsep=2pt]
  \item Completed table shown above.
  \item Completed pie chart shown above.
  \item White angle:
  \[
  \frac{20}{60}\times360=\frac13\times360=120^\circ.
  \]
  \item Black fraction \(=\frac{15}{60}=\frac14\).
  \item Angle sum check:
  \[
  120+90+72+54+24=360^\circ.
  \]
\end{enumerate}
\end{mdframed}

\vspace{6pt}

% ===========================
%  CHART 5
% ===========================
\begin{mdframed}[backgroundcolor=boxblue, linecolor=slice5,
                 linewidth=1.2pt, roundcorner=6pt]
\textbf{Chart 5 -- Languages Spoken at Home (survey of 90 students)}

\medskip
\begin{minipage}[c]{0.42\textwidth}
  \centering
  \piechart{4.0}{English/180/slice1,Urdu/72/slice2,Polish/36/slice3,Punjabi/36/slice4,Other/36/slice5}
\end{minipage}%
\hfill
\begin{minipage}[c]{0.54\textwidth}
  \centering
  \small
  \renewcommand{\arraystretch}{1.6}
  \begin{tabular}{|l|c|c|c|}
    \hline
    \textbf{Language} & \textbf{No.} & \textbf{Fraction} & \textbf{Angle} \\
    \hline
    English  & 45 & $\frac{1}{2}$ & $180^\circ$ \\
    \hline
    Urdu     & 18 & $\frac15$ & $72^\circ$ \\
    \hline
    Polish   &  9 & $\frac{1}{10}$ & $36^\circ$ \\
    \hline
    Punjabi  &  9 & $\frac{1}{10}$ & $36^\circ$ \\
    \hline
    Other    &  9 & $\frac{1}{10}$ & $36^\circ$ \\
    \hline
    \textbf{Total} & \textbf{90} & \textbf{1} & \textbf{360\textdegree} \\
    \hline
  \end{tabular}
\end{minipage}

\medskip
\textbf{Table working:}
\[
\frac{45}{90}=\frac12 \Rightarrow \frac12\times360=180^\circ,\qquad
\frac{18}{90}=\frac15 \Rightarrow \frac15\times360=72^\circ,
\]
\[
\frac{9}{90}=\frac{1}{10} \Rightarrow \frac{1}{10}\times360=36^\circ
\quad\text{for Polish, Punjabi and Other.}
\]
Check: \(180+72+36+36+36=360^\circ\).

\textbf{Answers:}
\begin{enumerate}[leftmargin=1.6em, itemsep=2pt]
  \item Completed table shown above.
  \item Completed pie chart shown above.
  \item Polish, Punjabi and Other all have 9 speakers, so each sector is
  \[
  \frac{9}{90}\times360=\frac{1}{10}\times360=36^\circ.
  \]
  \item English fraction \(=\frac{45}{90}=\frac12\).
  \item \textbf{Challenge:}
  \[
  \text{New total}=90+10=100,
  \qquad
  \text{New number of Polish speakers}=9+10=19.
  \]
  \[
  \text{New Polish angle}=\frac{19}{100}\times360=68.4^\circ.
  \]
  So the new sector angle for Polish is \(68.4^\circ\) (accept \(68^\circ\) or \(69^\circ\) to the nearest degree).
\end{enumerate}
\end{mdframed}

\end{document}